\section{Dual Scheduling Mechanism}

\subsection{Shapley Value Computation}

\textbf{Computational Challenge:} Exact Shapley value computation requires evaluating all $2^N$ possible coalitions, which is computationally prohibitive. We address this using MC-Shapley, a Monte Carlo permutation sampling approach combined with exponential smoothing.

\textbf{MC-Shapley Approximation:} The Shapley value for client $n$ in round $t$ is defined over all possible coalitions:
\begin{equation}
\varphi_n^{\text{current}}(t) = \sum_{S \subseteq \mathcal{S}_t \setminus \{n\}} \frac{|S|!(|\mathcal{S}_t| - |S| - 1)!}{|\mathcal{S}_t|!} \left[v(S \cup \{n\}) - v(S)\right]
\end{equation}
where $v(S)$ denotes the utility (negative validation loss) of the aggregated model from coalition $S$. This exact formulation requires enumerating all $2^{|\mathcal{S}_t|}$ coalitions, which is computationally prohibitive. To address this, our MC-Shapley method adopts \textit{Monte Carlo permutation sampling} as a practical approximation: in each iteration $m$, we randomly sample a permutation $\pi_m$ of the selected clients, then compute the marginal contribution of each client $n = \pi_m(j)$ as:
\begin{equation}
\Delta_n^{(m)} = v(\{\pi_m(1), \ldots, \pi_m(j)\}) - v(\{\pi_m(1), \ldots, \pi_m(j-1)\})
\end{equation}
The Shapley value is estimated as the average marginal contribution over $M$ random permutations:
\begin{equation}
\varphi_n^{\text{current}}(t) \approx \frac{1}{M} \sum_{m=1}^{M} \Delta_n^{(m)}
\end{equation}
This reduces the computational complexity from $O(2^K)$ to $O(K \times M)$ where $M \ll 2^K$. We further apply two truncation strategies to accelerate computation: (1) \textit{inter-round truncation}: if $|v(w_t) - v(w_{t-1})| < \epsilon$, skip Shapley computation; (2) \textit{intra-round truncation}: within a permutation, stop evaluating when $|v(S_j) - v(\mathcal{S}_t)| < \epsilon$.

\textbf{Critical Innovation — Exponential Smoothing:} A key design choice is the use of exponential smoothing to update Shapley values across rounds, balancing responsiveness to recent contributions with stability from historical observations:
\begin{equation}
\varphi_n(t) = \alpha \cdot \varphi_n(t-1) + (1 - \alpha) \cdot \varphi_n^{\text{current}}(t)
\end{equation}
where $\alpha \in [0, 1]$ is the smoothing parameter. A larger $\alpha$ gives more weight to historical values (noise filtering), while a smaller $\alpha$ prioritizes recent observations (fast adaptation). We set $\alpha = 0.5$ to balance noise filtering and rapid adaptation to changing client contributions in non-IID settings.

\subsection{Energy-Aware Scoring}

For each client $n$, we compute a normalized energy score:
\begin{equation}
s_{\text{energy},n}(t) = \frac{E_{\text{remain},n}(t)}{E_{\text{init}}}
\end{equation}

Before scoring, we filter out clients whose energy has fallen below the threshold:
\begin{equation}
\mathcal{A}(t) = \{n \mid E_{\text{remain},n}(t) \geq E_{\text{threshold}}\}
\end{equation}

At each round, we generate channel coefficients $h_n(t) \sim \mathcal{CN}(0, 1)$ for all clients and compute energy consumption as defined in Section III.

\subsection{Dual Scheduling Mechanism}

For each eligible client $n \in \mathcal{A}(t)$, we compute the composite score derived from Lyapunov drift-plus-penalty optimization:
\begin{equation}
\text{Score}_n(t) = V \cdot \hat{\varphi}_n(t) - \hat{Q}_n(t) \cdot \hat{E}_n(t)
\end{equation}
where $\hat{\varphi}_n(t)$ is the min-max normalized Shapley value, $\hat{E}_n(t)$ is the normalized energy consumption, and $\hat{Q}_n(t) = Q_n(t) / \max_{i} Q_i(t)$ is the normalized per-client virtual queue. Normalizing $Q_n(t)$ prevents the queue magnitude from overwhelming the Shapley value term as training progresses.

\textbf{Top-K Selection:} We employ a greedy selection strategy: (1) Sort all eligible clients by $\text{Score}_n(t)$ in descending order; (2) Select the top-$K$ clients: $\mathcal{S}_t = \text{top-}K(\mathcal{A}(t))$.

\subsection{Lyapunov-Based Dynamic Weight Adjustment}

To balance Shapley-based client quality and energy sustainability, we adopt Lyapunov optimization with per-client virtual queues. For each client $n$, we maintain a virtual queue $Q_n(t)$ that tracks cumulative energy budget violations.

After each round $t$, for selected clients, we update:
\begin{equation}
Q_n(t+1) = \max\left\{Q_n(t) + E_n(t) - \bar{E}, \; 0\right\}
\end{equation}
where $E_n(t)$ is the actual energy consumed by client $n$ in round $t$, and $\bar{E}$ is the per-round energy budget. When a client's energy consumption consistently exceeds the budget, its queue value $Q_n(t)$ grows, increasing the energy penalty in the scoring function and naturally steering selection away from energy-intensive clients. This per-client penalty mechanism adapts individually to each client's energy consumption pattern, providing finer-grained control than global weight adjustment.

\subsection{AES-256-GCM Gradient Encryption}

To protect gradient updates during transmission under the honest-but-curious server threat model, we integrate AES-256-GCM authenticated encryption as a transmission-layer security mechanism.

\textbf{Per-Client Key Management:} At system initialization, each client $n$ is assigned an independent 256-bit AES session key $k_n$. Keys are established once and reused across all rounds:
\begin{equation}
k_n \xleftarrow{\$} \{0,1\}^{256}, \quad \forall n \in \mathcal{N}
\end{equation}

\textbf{Encrypt-then-Upload:} After local training, client $n$ serializes its model update and encrypts it with a fresh nonce per round:
\begin{equation}
(c_n, \tau_n) = \text{AES-GCM}_{k_n}\!\left(\text{Serialize}(\Delta w_n),\; r_n\right)
\end{equation}
where $r_n \xleftarrow{\$} \{0,1\}^{96}$ is a fresh 96-bit nonce, $c_n$ is the ciphertext, and $\tau_n \in \{0,1\}^{128}$ is the GCM authentication tag.

\textbf{Server-Side Ephemeral Decryption:} The server decrypts, verifies the tag, uses the plaintext for Shapley computation and FedAvg aggregation, then immediately destroys the plaintext reference:
\begin{equation}
\Delta w_n = \text{AES-GCM-Dec}_{k_n}(c_n, \tau_n, r_n), \quad \text{then } \mathtt{del}(\Delta w_n)
\end{equation}
Only ciphertext $(c_n, \tau_n, r_n)$ is retained for audit purposes. This protocol ensures that no plaintext gradient persists on the server beyond the duration of a single round's computation.

\subsection{Complete Algorithm Framework}

\begin{algorithm}
\caption{Dual Scheduling with AES-256-GCM Secure Transmission}
\begin{algorithmic}[1]
\STATE \textbf{Input:} $N$ clients, $K$ selection size, $T$ rounds, $\alpha$ (smoothing), $V$ (Lyapunov), $\bar{E}$ (energy budget), $R$ (initial rounds), $\{k_n\}$ (AES-256 session keys)
\STATE \textbf{Output:} Global model $w_T$
\STATE Initialize: $w_0$, $E_{\text{remain},n}(0) = E_{\text{init}}$, $\varphi_n(0) = 0$, $Q_n(0) = 0$, $\forall n$
\FOR{$t = 1$ to $T$}
    \STATE Compute available clients: $\mathcal{A}(t) = \{n \mid E_{\text{remain},n}(t) \geq E_{\text{threshold}}\}$
    \IF{$t \leq R$}
        \STATE $\mathcal{S}_t \leftarrow$ round-robin selection of $K$ clients from $\mathcal{A}(t)$
    \ELSE
        \STATE Normalize: $\hat{\varphi}_n(t) = (\varphi_n(t) - \varphi_{\min}) / (\varphi_{\max} - \varphi_{\min})$, $\hat{E}_n(t) = E_n(t) / E_{\max}(t)$, $\hat{Q}_n(t) = Q_n(t) / \max_i Q_i(t)$
        \FOR{each $n \in \mathcal{A}(t)$}
            \STATE $\text{Score}_n(t) = V \cdot \hat{\varphi}_n(t) - \hat{Q}_n(t) \cdot \hat{E}_n(t)$
        \ENDFOR
        \STATE $\mathcal{S}_t \leftarrow \text{top-}K$ clients from $\mathcal{A}(t)$ by $\text{Score}_n(t)$
    \ENDIF
    \STATE Broadcast $w_t$ to clients in $\mathcal{S}_t$
    \FOR{each client $n \in \mathcal{S}_t$ (in parallel)}
        \STATE Perform local training for $E$ epochs on $\mathcal{D}_n$
        \STATE Compute update: $\Delta w_n = w_n^{\text{local}} - w_t$
        \STATE Generate nonce $r_n \xleftarrow{\$} \{0,1\}^{96}$
        \STATE Encrypt: $(c_n, \tau_n) = \text{AES-GCM}_{k_n}(\text{Serialize}(\Delta w_n), r_n)$
        \STATE Upload $(c_n, \tau_n, r_n)$ to server
    \ENDFOR
    \FOR{each $n \in \mathcal{S}_t$}
        \STATE Decrypt: $\Delta w_n \leftarrow \text{AES-GCM-Dec}_{k_n}(c_n, \tau_n, r_n)$
    \ENDFOR
    \STATE $w_{t+1} = \sum_{n\in\mathcal{S}_t} \frac{|\mathcal{D}_n|}{\sum_{i\in\mathcal{S}_t}|\mathcal{D}_i|} w_n^{\text{local}}$
    \STATE \textbf{Destroy} all plaintext $\Delta w_n$ references
    \STATE Update Shapley values using MC-Shapley with exponential smoothing
    \FOR{each $n \in \mathcal{S}_t$}
        \STATE Generate channel: $h_n(t) \sim \mathcal{CN}(0,1)$
        \STATE $E_n(t) = E_{\text{trans},n}(t) + E_{\text{comp},n}(t)$
        \STATE $E_{\text{remain},n}(t+1) = E_{\text{remain},n}(t) - E_n(t)$
        \STATE $Q_n(t+1) = \max\{Q_n(t) + E_n(t) - \bar{E}, \; 0\}$
    \ENDFOR
\ENDFOR
\STATE \textbf{return} $w_T$
\end{algorithmic}
\end{algorithm}
